\begingroup
\setlength{\abovedisplayskip}{6pt}
\setlength{\belowdisplayskip}{6pt}
% Printed82; principalPDF89. Clipped final column supplied by HKUST page00090.
\clearpage
\noindent\makebox[12mm][l]{82}\makebox[\dimexpr\linewidth-24mm\relax][c]{\textsc{Nachlass.}}\makebox[12mm][r]{}\par\medskip
\begin{center}
\begingroup\fontsize{7.5}{9.0}\selectfont
\setlength{\tabcolsep}{1.8pt}\renewcommand{\arraystretch}{1.03}
\begin{tabular}{r|r|r|r|r|r|r|r|r|r}\hline
\(51^\circ\)&\shortstack{\([\log A'\)\\\(8{,}510\ldotp\ldotp\ldotp\ldotp-10\)}&
\shortstack{\(\log B'\)\\\(8{,}5089\ldotp\ldotp\ldotp-10\)}&
\shortstack{\(\log10^7.F'\)\\\(1{,}9505.-10\)}&
\shortstack{\(\log10^7.G'\)\\\(9{,}5\ldotp\ldotp\ldotp\ldotp-20\)}&
\shortstack{\(\log10^7.D'\)\\\(4{,}61\ldotp\ldotp\ldotp-10\)}&
\shortstack{\(\log10^7.I'\)\\\(3{,}24\ldotp\ldotp\ldotp-10\)}&
\shortstack{\(\log(-10^7.H')\)\\\(2{,}2\ldotp\ldotp\ldotp\ldotp-10\)}&
\shortstack{\(\log10^7.C'\)\\\(4{,}626\ldotp\ldotp-10\)}&
\shortstack{\(\log10^7.E'\)\\\(4{,}9308.-10]\)}\\\hline\hline
\(0'\) & \(8{,}510\,1012\) & \(8{,}508\,9610\) & \(1{,}95055\) & \(9{,}55074\) & \(4{,}61525\) & \(3{,}24181\) & \(2{,}22778\) & \(4{,}62643\) & \(4{,}93089\) \\
\(1\) & \(8{,}510\,1000\) & \(05\) &  & \(094\) & \(25\) & \(85\) & \(2{,}22902\) & \(4{,}62643\) &  \\
\(2\) & \(8{,}510\,0988\) & \(8{,}508\,9601\) &  & \(115\) & \(24\) & \(89\) & \(2{,}23025\) & \(4{,}62644\) &  \\
\(3\) & \(76\) & \(8{,}508\,9597\) & \(1{,}95055\) & \(135\) & \(24\) & \(93\) & \(148\) &  &  \\
\(4\) & \(63\) & \(93\) & \(1{,}95054\) & \(155\) & \(23\) & \(3{,}24197\) & \(271\) &  &  \\
\(5\) & \(8{,}510\,0951\) & \(8{,}508\,9589\) &  & \(9{,}55176\) & \(4{,}61523\) & \(3{,}24201\) & \(2{,}23394\) &  &  \\
\hline
\(6\) & \(8{,}510\,0939\) & \(8{,}508\,9585\) &  & \(9{,}55196\) & \(4{,}61522\) & \(3{,}24206\) & \(2{,}23516\) &  &  \\
\(7\) & \(27\) & \(81\) &  & \(216\) & \(22\) & \(10\) & \(637\) & \(4{,}62644\) &  \\
\(8\) & \(14\) & \(77\) &  & \(236\) & \(21\) & \(14\) & \(759\) & \(4{,}62645\) &  \\
\(9\) & \(8{,}510\,0902\) & \(73\) &  & \(257\) & \(21\) & \(18\) & \(2{,}23880\) &  & \(4{,}93089\) \\
\(10\) & \(8{,}510\,0890\) & \(8{,}508\,9569\) &  & \(9{,}55277\) & \(4{,}61520\) & \(3{,}24222\) & \(2{,}24000\) &  & \(4{,}93088\) \\
\hline
\(11\) & \(8{,}510\,0878\) & \(8{,}508\,9565\) &  & \(9{,}55297\) & \(4{,}61520\) & \(3{,}24226\) & \(2{,}24121\) &  &  \\
\(12\) & \(65\) & \(60\) &  & \(318\) & \(19\) & \(30\) & \(241\) &  &  \\
\(13\) & \(53\) & \(56\) &  & \(338\) & \(19\) & \(34\) & \(360\) & \(4{,}62645\) &  \\
\(14\) & \(41\) & \(52\) &  & \(358\) & \(18\) & \(38\) & \(480\) & \(4{,}62646\) &  \\
\(15\) & \(8{,}510\,0829\) & \(8{,}508\,9548\) &  & \(9{,}55378\) & \(4{,}61518\) & \(3{,}24242\) & \(2{,}24599\) &  &  \\
\hline
\(16\) & \(8{,}510\,0816\) & \(8{,}508\,9544\) & \(1{,}95054\) & \(9{,}55398\) & \(4{,}61517\) & \(3{,}24246\) & \(2{,}24717\) &  &  \\
\(17\) & \(8{,}510\,0804\) & \(40\) & \(1{,}95053\) & \(419\) & \(17\) & \(50\) & \(836\) &  &  \\
\(18\) & \(8{,}510\,0792\) & \(36\) &  & \(439\) & \(16\) & \(55\) & \(2{,}24953\) &  &  \\
\(19\) & \(80\) & \(32\) &  & \(459\) & \(16\) & \(59\) & \(2{,}25071\) & \(4{,}62646\) &  \\
\(20\) & \(8{,}510\,0767\) & \(8{,}508\,9528\) &  & \(9{,}55479\) & \(4{,}61515\) & \(3{,}24263\) & \(2{,}25188\) & \(4{,}62647\) &  \\
\hline
\(21\) & \(8{,}510\,0755\) & \(8{,}508\,9524\) &  & \(9{,}55499\) & \(4{,}61515\) & \(3{,}24267\) & \(2{,}25305\) &  &  \\
\(22\) & \(43\) & \(20\) &  & \(519\) & \(14\) & \(71\) & \(422\) &  & \(4{,}93088\) \\
\(23\) & \(31\) & \(16\) &  & \(539\) & \(14\) & \(75\) & \(538\) &  & \(4{,}93087\) \\
\(24\) & \(18\) & \(11\) &  & \(559\) & \(13\) & \(79\) & \(654\) &  &  \\
\(25\) & \(8{,}510\,0706\) & \(8{,}508\,9507\) &  & \(9{,}55580\) & \(4{,}61513\) & \(3{,}24283\) & \(2{,}25770\) & \(4{,}62647\) &  \\
\hline
\(26\) & \(8{,}510\,0694\) & \(8{,}508\,9503\) &  & \(9{,}55600\) & \(4{,}61512\) & \(3{,}24287\) & \(2{,}25885\) & \(4{,}62648\) &  \\
\(27\) & \(81\) & \(8{,}508\,9499\) &  & \(620\) & \(12\) & \(91\) & \(2{,}26000\) &  &  \\
\(28\) & \(69\) & \(95\) & \(1{,}95053\) & \(640\) & \(11\) & \(95\) & \(115\) &  &  \\
\(29\) & \(57\) & \(91\) & \(1{,}95052\) & \(660\) & \(11\) & \(3{,}24299\) & \(229\) &  &  \\
\(30\) & \(8{,}510\,0645\) & \(8{,}508\,9487\) &  & \(9{,}55680\) & \(4{,}61510\) & \(3{,}24303\) & \(2{,}26343\) &  &  \\
\hline
\(31\) & \(8{,}510\,0633\) & \(8{,}508\,9483\) &  & \(9{,}55700\) & \(4{,}61509\) & \(3{,}24307\) & \(2{,}26457\) & \(4{,}62648\) &  \\
\(32\) & \(21\) & \(79\) &  & \(720\) & \(09\) & \(11\) & \(570\) & \(4{,}62649\) &  \\
\(33\) & \(8{,}510\,0608\) & \(75\) &  & \(740\) & \(08\) & \(16\) & \(684\) &  &  \\
\(34\) & \(8{,}510\,0596\) & \(71\) &  & \(760\) & \(08\) & \(20\) & \(796\) &  & \(4{,}93087\) \\
\(35\) & \(8{,}510\,0584\) & \(8{,}508\,9467\) &  & \(9{,}55780\) & \(4{,}61507\) & \(3{,}24324\) & \(2{,}26909\) &  & \(4{,}93086\) \\
\hline
\(36\) & \(8{,}510\,0572\) & \(8{,}508\,9463\) &  & \(9{,}55800\) & \(4{,}61507\) & \(3{,}24328\) & \(2{,}27021\) &  &  \\
\(37\) & \(59\) & \(59\) &  & \(820\) & \(06\) & \(32\) & \(133\) & \(4{,}62649\) &  \\
\(38\) & \(47\) & \(54\) &  & \(839\) & \(06\) & \(36\) & \(245\) & \(4{,}62650\) &  \\
\(39\) & \(35\) & \(50\) &  & \(859\) & \(05\) & \(40\) & \(356\) &  &  \\
\(40\) & \(8{,}510\,0523\) & \(8{,}508\,9446\) & \(1{,}95052\) & \(9{,}55879\) & \(4{,}61505\) & \(3{,}24344\) & \(2{,}27467\) &  &  \\
\hline
\(41\) & \(8{,}510\,0511\) & \(8{,}508\,9442\) & \(1{,}95051\) & \(9{,}55899\) & \(4{,}61504\) & \(3{,}24348\) & \(2{,}27578\) &  &  \\
\(42\) & \(8{,}510\,0498\) & \(38\) &  & \(919\) & \(04\) & \(52\) & \(688\) &  &  \\
\(43\) & \(86\) & \(34\) &  & \(939\) & \(03\) & \(56\) & \(798\) & \(4{,}62650\) &  \\
\(44\) & \(74\) & \(30\) &  & \(959\) & \(03\) & \(60\) & \(2{,}27908\) & \(4{,}62651\) &  \\
\(45\) & \(8{,}510\,0462\) & \(8{,}508\,9426\) &  & \(9{,}55979\) & \(4{,}61502\) & \(3{,}24364\) & \(2{,}28018\) &  &  \\
\hline
\(46\) & \(8{,}510\,0450\) & \(8{,}508\,9422\) &  & \(9{,}55999\) & \(4{,}61502\) & \(3{,}24368\) & \(2{,}28127\) &  & \(4{,}93086\) \\
\(47\) & \(37\) & \(18\) &  & \(9{,}56018\) & \(01\) & \(72\) & \(236\) &  & \(4{,}93085\) \\
\(48\) & \(25\) & \(14\) &  & \(038\) & \(01\) & \(76\) & \(344\) &  &  \\
\(49\) & \(13\) & \(10\) &  & \(058\) & \(00\) & \(80\) & \(453\) & \(4{,}62651\) &  \\
\(50\) & \(8{,}510\,0401\) & \(8{,}508\,9406\) &  & \(9{,}56078\) & \(4{,}61500\) & \(3{,}24384\) & \(2{,}28561\) & \(4{,}62652\) &  \\
\hline
\(51\) & \(8{,}510\,0389\) & \(8{,}508\,9402\) &  & \(9{,}56098\) & \(4{,}61499\) & \(3{,}24388\) & \(2{,}28669\) &  &  \\
\(52\) & \(76\) & \(8{,}508\,9398\) &  & \(117\) & \(99\) & \(93\) & \(776\) &  &  \\
\(53\) & \(64\) & \(93\) & \(1{,}95051\) & \(137\) & \(98\) & \(3{,}24397\) & \(884\) &  &  \\
\(54\) & \(52\) & \(89\) & \(1{,}95050\) & \(157\) & \(98\) & \(3{,}24401\) & \(2{,}28990\) &  &  \\
\(55\) & \(8{,}510\,0340\) & \(8{,}508\,9385\) &  & \(9{,}56176\) & \(4{,}61497\) & \(3{,}24405\) & \(2{,}29097\) & \(4{,}62652\) &  \\
\hline
\(56\) & \(8{,}510\,0328\) & \(8{,}508\,9381\) &  & \(9{,}56196\) & \(4{,}61497\) & \(3{,}24409\) & \(2{,}29204\) & \(4{,}62653\) &  \\
\(57\) & \(15\) & \(77\) &  & \(216\) & \(96\) & \(13\) & \(310\) &  &  \\
\(58\) & \(8{,}510\,0303\) & \(73\) &  & \(236\) & \(96\) & \(17\) & \(416\) &  &  \\
\(59\) & \(8{,}510\,0291\) & \(69\) &  & \(255\) & \(95\) & \(21\) & \(521\) &  & \(4{,}93085\) \\
\(60\) & \(8{,}510\,0279\) & \(8{,}508\,9365\) & \(1{,}95050\) & \(9{,}56275\) & \(4{,}61495\) & \(3{,}24425\) & \(2{,}29627\) & \(4{,}62653\) & \(4{,}93084\) \\
\end{tabular}
\endgroup
\end{center}
\noindent[Bei der Berechnung dieser Tabelle wurde der Abplattungswert \(\frac1{302{,}78}\) verwendet.]

% Printed83; principalPDF90.
\clearpage
\noindent\makebox[12mm][l]{}\makebox[\dimexpr\linewidth-24mm\relax][c]{\textsc{Geodätische Linie.}}\makebox[12mm][r]{83}\par\medskip
\begin{center}[5.]\end{center}
Das Ergebnis der neuen Methode für die kürzeste Linie auf dem Sphäroid
lautet [wenn das Azimut von Süden nach Westen und die Länge nach Osten positiv gezählt werden]:
\begingroup\small
\[
\begin{array}{@{}lcl@{}}
\displaystyle A=\frac{206\,265(1-ee\sin\varphi^2)^{\frac32}}{a(1-ee)}
&\text{[Logarithmus:]}&8{,}512\,7017.130-10-\tfrac32\mathrm{C}.\log(1-ee\sin\varphi^2)\\[2ex]
\displaystyle B=\frac{206\,265(1-ee\sin\varphi^2)^{\frac12}}a
&\text{»}&8{,}509\,8272.984-10-\tfrac12\mathrm{C}.\log(1-ee\sin\varphi^2)\\[2ex]
\displaystyle\alpha=\frac{\mathrm{Mod}.10^7}{24.206\,265^2}
&\text{»}&4{,}628\,7228.053-10\\[2ex]
\displaystyle\beta=\alpha.\frac{1-ee\sin\varphi^2}{1-ee}
&\text{»}&4{,}631\,5972.199-10-\mathrm{C}.\log(1-ee\sin\varphi^2)\\[2ex]
\displaystyle\gamma=\frac\beta{aa}.206\,265^2
&\text{»}&1{,}651\,2518.167-10-\mathrm{C}.\log(1-ee\sin\varphi^2)\\[2ex]
\displaystyle\delta=\frac{\alpha(1-ee)}{(1-ee\sin\varphi^2)^2}
&\text{»}&4{,}625\,8483.907-10+2\mathrm{C}.\log(1-ee\sin\varphi^2).
\end{array}
\]
\endgroup
\noindent[\(\alpha,\beta,\gamma,\delta\) werden in Einheiten der 7. Dezimalstelle angegeben.
\(A\) und \(B\) wurden im vorherigen Artikel mit \(A'\) und \(B'\) bezeichnet.
Bei den Zahlenwerten wurden die Abplattung \(\frac1{302{,}68}\) und der entsprechende Wert
\[
\log(1-ee)=9{,}997\,1255.854
\]
\noindent verwendet.]
\[
\begin{aligned}
\log a&=6{,}804\,5978.348\\
\log206\,264{,}806\ldotp\ldotp&=5{,}314\,4251.332.
\end{aligned}
\]
\[
\begin{array}{rcl@{\qquad}l}
\mathrm I&=&\alpha.\delta\zeta^2&\log(1)=\mathrm I+2\mathrm{II}\\
\mathrm{II}&=&\beta.\delta\lambda^2&\log(2)=\mathrm I+2\mathrm{III}\\
\mathrm{III}&=&\gamma.ss&\log(3)=\mathrm{II}-\mathrm{III}=\mathrm I-\mathrm{IV}\\
\mathrm{IV}&=&\delta.\delta\varphi^2&\text{Näherungsweise }\mathrm I+\mathrm{III}=\mathrm{II}+\mathrm{IV}.
\end{array}
\]
\[
\begin{aligned}
\Delta\varphi&=-(1)A.s\cos\zeta\\
\Delta\zeta&=-(2)B.s\sin\zeta\operatorname{tang}\varphi\\
\Delta\lambda&=-\frac{(3)B.s\sin\zeta}{\cos\varphi}.
\end{aligned}
\]
Bei diesen Rechnungen sind für \(\zeta\) und \(\varphi\) die Mittelwerte der an den
beiden Punkten geltenden Werte zugrunde zu legen. [Die Koeffizienten
\(A,B,\alpha,\beta,\gamma,\delta\) gehören zu diesem mittleren \(\varphi\). Die Größe
\(\varepsilon=\rho.\frac{AB}{2}=\frac{206265(1-ee\sin\varphi^2)^2}{2aa(1-ee)}\)
in der Tabelle auf S.~84 dient zur Berechnung des Exzesses.]
\par\smallskip\noindent\hfill\(11^*\)\hspace{.12\textwidth}\par

% Printed84; principalPDF91; GooglePDF95 only for isolated-mark comparison.
\clearpage
\noindent\makebox[12mm][l]{84}\makebox[\dimexpr\linewidth-24mm\relax][c]{\textsc{Nachlass.}}\makebox[12mm][r]{}\par\medskip
\begin{center}
\begingroup\fontsize{8}{9.3}\selectfont
\setlength{\tabcolsep}{5pt}\renewcommand{\arraystretch}{1.03}
\begin{tabular}{r||r|r|r|r|r|r}\hline
\(50^\circ\)&\shortstack{\([\log A\)\\\(8{,}5101\ldotp\ldotp\ldotp-10\)}&
\shortstack{\(\log B\)\\\(8{,}5089\ldotp\ldotp\ldotp-10\)}&
\shortstack{\(\log\beta\)\\\(4{,}629\ldotp\ldotp-10\)}&
\shortstack{\(\log\gamma\)\\\(1{,}6495.-10\)}&
\shortstack{\(\log\delta\)\\\(4{,}629\ldotp\ldotp-10\)}&
\shortstack{\(\log\varepsilon\)\\\(1{,}403\ldotp\ldotp-10]\)}\\\hline\hline
\(0'\) & \(8{,}510\,1751\) & \(8{,}508\,9856\) & \(4{,}62991\) & \(1{,}64957\) & \(4{,}62922\) & \(1{,}40371\) \\
\(1\) & \(39\) & \(52\) &  &  &  & \(1{,}40370\) \\
\(2\) & \(27\) & \(48\) &  & \(1{,}64957\) &  &  \\
\(3\) & \(14\) & \(43\) &  & \(1{,}64956\) &  &  \\
\(4\) & \(8{,}510\,1702\) & \(39\) &  &  & \(4{,}62922\) &  \\
\(5\) & \(8{,}510\,1690\) & \(8{,}508\,9835\) &  &  & \(4{,}62923\) &  \\
\hline
\(6\) & \(77\) & \(8{,}508\,9831\) &  &  &  & \(1{,}40370\) \\
\(7\) & \(65\) & \(27\) &  &  &  & \(1{,}40369\) \\
\(8\) & \(53\) & \(23\) & \(4{,}62991\) &  &  &  \\
\(9\) & \(40\) & \(19\) & \(4{,}62990\) &  &  &  \\
\(10\) & \(28\) & \(8{,}508\,9815\) &  &  & \(4{,}62923\) &  \\
\hline
\(11\) & \(15\) & \(8{,}508\,9811\) &  &  & \(4{,}62924\) &  \\
\(12\) & \(8{,}510\,1603\) & \(06\) &  &  &  & \(1{,}40369\) \\
\(13\) & \(8{,}510\,1591\) & \(8{,}508\,9802\) &  &  &  & \(1{,}40368\) \\
\(14\) & \(78\) & \(8{,}508\,9798\) &  &  &  &  \\
\(15\) & \(66\) & \(8{,}508\,9794\) &  & \(1{,}64956\) &  &  \\
\hline
\(16\) & \(54\) & \(8{,}508\,9790\) &  & \(1{,}64955\) & \(4{,}62924\) &  \\
\(17\) & \(41\) & \(86\) &  &  & \(4{,}62925\) &  \\
\(18\) & \(29\) & \(82\) &  & \raisebox{.6ex}{.} &  & \(1{,}40368\) \\
\(19\) & \(17\) & \(78\) &  &  &  & \(1{,}40367\) \\
\(20\) & \(8{,}510\,1504\) & \(8{,}508\,9774\) &  &  &  &  \\
\hline
\(21\) & \(8{,}510\,1492\) & \(8{,}508\,9769\) & \(4{,}62990\) &  &  &  \\
\(22\) & \(80\) & \(65\) & \(4{,}62989\) &  & \(4{,}62925\) &  \\
\(23\) & \(67\) & \(61\) &  &  & \(4{,}62926\) &  \\
\(24\) & \(55\) & \(57\) &  &  &  & \(1{,}40367\) \\
\(25\) & \(43\) & \(8{,}508\,9753\) &  &  &  & \(1{,}40366\) \\
\hline
\(26\) & \(31\) & \(8{,}508\,9749\) &  &  &  &  \\
\(27\) & \(18\) & \(45\) &  & \(1{,}64955\) &  &  \\
\(28\) & \(8{,}510\,1406\) & \(41\) &  & \(1{,}64954\) &  &  \\
\(29\) & \(8{,}510\,1394\) & \(37\) &  &  & \(4{,}62926\) &  \\
\(30\) & \(81\) & \(8{,}508\,9732\) &  &  & \(4{,}62927\) & \(1{,}40366\) \\
\hline
\(31\) & \(69\) & \(8{,}508\,9728\) &  &  &  & \(1{,}40365\) \\
\(32\) & \(57\) & \(24\) &  &  &  &  \\
\(33\) & \(44\) & \(20\) & \(4{,}62989\) &  &  &  \\
\(34\) & \(32\) & \(16\) & \(4{,}62988\) &  &  &  \\
\(35\) & \(20\) & \(8{,}508\,9712\) &  &  & \(4{,}62927\) &  \\
\hline
\(36\) & \(8{,}510\,1307\) & \(8{,}508\,9708\) &  &  & \(4{,}62928\) & \(1{,}40365\) \\
\(37\) & \(8{,}510\,1295\) & \(04\) &  &  &  & \(1{,}40364\) \\
\(38\) & \(83\) & \(8{,}508\,9700\) &  &  &  &  \\
\(39\) & \(70\) & \(8{,}508\,9696\) &  & \(1{,}64954\) &  &  \\
\(40\) & \(58\) & \(8{,}508\,9691\) &  & \(1{,}64953\) &  &  \\
\hline
\(41\) & \(46\) & \(8{,}508\,9687\) &  &  & \(4{,}62928\) &  \\
\(42\) & \(34\) & \(83\) &  &  & \(4{,}62929\) & \(1{,}40364\) \\
\(43\) & \(21\) & \(79\) &  &  &  & \(1{,}40363\) \\
\(44\) & \(8{,}510\,1209\) & \(75\) &  &  &  &  \\
\(45\) & \(8{,}510\,1197\) & \(8{,}508\,9671\) & \(4{,}62988\) &  &  &  \\
\hline
\(46\) & \(84\) & \(8{,}508\,9667\) & \(4{,}62987\) &  &  &  \\
\(47\) & \(72\) & \(63\) &  &  & \(4{,}62929\) &  \\
\(48\) & \(60\) & \(59\) &  &  & \(4{,}62930\) &  \\
\(49\) & \(48\) & \(55\) &  &  &  & \(1{,}40363\) \\
\(50\) & \(35\) & \(8{,}508\,9650\) &  &  &  & \(1{,}40362\) \\
\hline
\(51\) & \(23\) & \(8{,}508\,9646\) &  & \(1{,}64953\) &  &  \\
\(52\) & \(8{,}510\,1111\) & \(42\) &  & \(1{,}64952\) &  &  \\
\(53\) & \(8{,}510\,1098\) & \(38\) &  &  & \(4{,}62930\) &  \\
\(54\) & \(86\) & \(34\) &  &  & \(4{,}62931\) &  \\
\(55\) & \(74\) & \(8{,}508\,9630\) &  &  &  & \(1{,}40362\) \\
\hline
\(56\) & \(62\) & \(8{,}508\,9626\) &  &  &  & \(1{,}40361\) \\
\(57\) & \(49\) & \(22\) & \(4{,}62987\) &  &  &  \\
\(58\) & \(37\) & \(18\) & \(4{,}62986\) &  & \(4{,}62931\) &  \\
\(59\) & \(25\) & \(14\) &  &  & \(4{,}62932\) &  \\
\(60\) & \(8{,}510\,1012\) & \(8{,}508\,9610\) & \(4{,}62986\) & \(1{,}64952\) & \(4{,}62932\) & \(1{,}40361\) \\
\end{tabular}
\endgroup
\end{center}
\noindent[Bei der Berechnung dieser Tabelle wurde der Abplattungswert \(\frac1{302{,}78}\) verwendet.]

% Printed85; principalPDF92.
\clearpage
\noindent\makebox[12mm][l]{}\makebox[\dimexpr\linewidth-24mm\relax][c]{\textsc{Geodätische Linie.}}\makebox[12mm][r]{85}\par\medskip
\begin{center}[6.]\end{center}
Das Vorstehende [Art.~4 und 5] ist die indirekte Lösung. Die direkte Lösung
erhält man folgendermaßen:
\[
\begin{aligned}
\surd(1-ee).\operatorname{tang}\varphi&=\operatorname{tang}\psi\\
\surd(1-ee\sin\varphi^2)&=\frac{\cos\varphi}{\cos\psi}
=\sqrt{\frac{1-ee}{1-ee\cos\psi^2}}\\
S&=\frac sa\sqrt{\frac{1-ee\sin\varphi^2}{1-ee}}
\end{aligned}
\]
\noindent[wobei \(s\) die lineare Länge der kürzesten Linie auf dem Ellipsoid ist].
(Eigentlich sollte man hier anstelle von \(\varphi\) den Wert \(\tfrac12(\varphi+\varphi')\) verwenden).
Man löse das sphärische Dreieck mit den nachstehenden Seiten und Winkeln:
\[
\begin{array}{c|c}
\text{Seiten}&\text{Winkel}\\\hline
90^\circ-\psi&\zeta'\\
90^\circ-\psi'&180^\circ-\zeta\\
S&L
\end{array}
\]
\noindent[Vgl. Band IV, S.~286 ff.] Meine Formeln ergeben hier
\[
\begin{aligned}
\sin\tfrac12\zeta\sin(45^\circ-\tfrac12\psi-\tfrac12S)
&=\sin\tfrac12(\zeta'-L)\sin(45^\circ-\tfrac12\psi')\\
\cos\tfrac12\zeta\sin(45^\circ-\tfrac12\psi+\tfrac12S)
&=\cos\tfrac12(\zeta'-L)\sin(45^\circ+\tfrac12\psi')\\
\sin\tfrac12\zeta\cos(45^\circ-\tfrac12\psi-\tfrac12S)
&=\sin\tfrac12(\zeta'+L)\cos(45^\circ-\tfrac12\psi')\\
\cos\tfrac12\zeta\cos(45^\circ-\tfrac12\psi+\tfrac12S)
&=\cos\tfrac12(\zeta'+L)\cos(45^\circ+\tfrac12\psi'),
\end{aligned}
\]
\noindent und dann gilt
\[
\operatorname{tang}\varphi'=\frac{\operatorname{tang}\psi'}{\surd(1-ee)}
\qquad\text{[und näherungsweise]}\qquad
\begin{aligned}
\varphi-\varphi'&=(\psi-\psi')\frac{1-ee\sin\varphi^{*2}}{\surd(1-ee)}\\
&=(\psi-\psi')\frac{\surd(1-ee)}{1-ee\cos\psi^{*2}}
\end{aligned}
\]
\[
\lambda=L\surd(1-ee\cos\psi^{*2})
=L\sqrt{\frac{1-ee}{1-ee\sin\varphi^{*2}}}
\]
\[
[\varphi^*=\tfrac12(\varphi+\varphi'),\qquad
\psi^*=\tfrac12(\psi+\psi').
\]
\noindent\(\zeta\) wird von Süden nach Westen und \(\zeta'\) von Norden nach Osten gezählt. \(\lambda\) und \(L\) sind nach Westen positiv].

Vorteilhafter sind hier jedoch die folgenden Formeln:
\[
\begin{aligned}
\sin\psi'&=\sin\psi\cos S-\cos\psi\sin S\cos\zeta\\
\cos\psi'\cos\zeta'&=\sin\psi\sin S+\cos\psi\cos S\cos\zeta\\
\cos\psi'\sin\zeta'&=\cos\psi\sin\zeta.
\end{aligned}
\]

% Printed86; principalPDF93; unclear delimiters in3) follow HKUSTpg00094.
\clearpage
\noindent\makebox[12mm][l]{86}\makebox[\dimexpr\linewidth-24mm\relax][c]{\textsc{Nachlass.}}\makebox[12mm][r]{}\par\medskip
Also
\[
\operatorname{cotang}\zeta'=\cos S\operatorname{cotang}\zeta
+\frac{\surd(1-ee).\operatorname{tang}\varphi\sin S}{\sin\zeta}
\]
\noindent oder näherungsweise:
\[
\operatorname{cotang}\zeta'=\operatorname{cotang}\zeta
\left(1-\frac{ss}{2aa}.\frac{1-ee\sin\varphi^2}{1-ee}\right)
\left(1+\frac{\operatorname{tang}\varphi}{\cos\zeta}.\frac sa\surd(1-ee\sin\varphi^2)\right).
\]
\noindent Ferner gilt
\[
\begin{aligned}
\cos\psi'\cos L&=\cos\psi\cos S+\sin\psi\sin S\cos\zeta\\
\cos\psi'\sin L&=\sin S\sin\zeta.
\end{aligned}
\]
\noindent Näherungsweise bis zur \(2^{\text{ten}}\) Ordnung
\[
L=\frac{S\sin\zeta}{\cos\varphi}
-\frac{SS\operatorname{tang}\psi}{\cos\psi}\cos\zeta\sin\zeta.
\]
Das obige Problem lässt sich auch so lösen:
\[
\begin{array}{r@{\qquad}l}
1)&\operatorname{tang}\psi=\surd(1-ee).\operatorname{tang}\varphi\\[1ex]
2)&\operatorname{tang}u=\operatorname{tang}S\cos\zeta\\[1ex]
3)&\displaystyle\operatorname{tang}L=\frac{\sin u\operatorname{tang}\zeta}{\cos(\psi-u)}
=\frac{\operatorname{tang}S\sin p}{\cos(\psi-u)}\\[1.5ex]
4)&\operatorname{tang}\psi'=\operatorname{tang}(\psi-u)\cos L\\[1ex]
5)&\displaystyle\operatorname{tang}\varphi'=\frac{\operatorname{tang}\psi'}{\surd(1-ee)}\\[1.5ex]
6)&\operatorname{tang}p=\operatorname{tang}\zeta\cos S\\[1ex]
7)&\operatorname{tang}q=\operatorname{tang}(\psi-u)\sin\zeta\sin S=\operatorname{tang}L\sin\psi'\\[1ex]
8)&p-q=\zeta'\\[1ex]
9)&\displaystyle\lambda=L\sqrt{\frac{1-ee}{1-ee\sin\varphi^2}}.
\end{array}
\]
Setzt man
\[
\operatorname{tang}(\psi-u)=\surd(1-ee).\operatorname{tang}(\varphi-w),
\]
\noindent so gilt mit hinreichender Genauigkeit
\[
w=(4)A.s\cos\zeta,
\]
\begingroup\small
\[
\left[A=\frac{(1-ee\sin\varphi^2)^{\frac32}}{a(1-ee)}\quad\text{und}\right]
\quad\begin{aligned}
\log(4)&=\tfrac13\frac{ss}{aa}.\frac{1-ee\sin\varphi^2}{1-ee}
-\tfrac13ww\frac{1-ee}{(1-ee\sin\varphi^2)^2}\\
&=\mu-\nu.
\end{aligned}
\]
\endgroup

% Printed87; principalPDF94.
\clearpage
\noindent\makebox[12mm][l]{}\makebox[\dimexpr\linewidth-24mm\relax][c]{\textsc{Geodätische Linie.}}\makebox[12mm][r]{87}\par\medskip
\noindent[Ferner gilt näherungsweise]
\[
p=\zeta-\tfrac12(5).\frac{1-ee\sin\varphi^2}{1-ee}\frac{ss}{aa}\cos\zeta\sin\zeta
\]
\[\log(5)=\tfrac54\mu-\tfrac32\nu.\]
\begin{center}[7.]\\[.5ex]\textbf{Musterrechnung.}\end{center}
\noindent[Übertragung der geografischen Koordinaten von Mannheim nach Seeberg.]

[Hierbei wurde der Abplattungswert \(\frac1{302,68}\) verwendet.

1. Anwendung. der Formeln aus Art.~5.]
\begingroup\small
\[
\begin{array}{lrlrl}
\text{Mannheim}&\varphi=&49^\circ29'12'',930&\zeta=224^\circ18'2'',100&\log s=5,358\,6482\\
\text{Seeberg}&\varphi'=&50\quad56\quad5,514&\zeta'=226\quad2\quad41,480&\Delta\lambda=2^\circ16'10'',868\\\cline{2-4}
&\varphi^*=&50\quad12\quad39,222&\zeta^*=225\quad10\quad21,790&
\end{array}
\]
\noindent\begin{minipage}[t]{.28\linewidth}\vspace{0pt}\centering
\(\begin{array}{l@{\;}r}
s\ldotp\,\ldotp\,\ldotp\,\ldotp\,\ldotp\,\ldotp&5,358\,6482\\
A\ldotp\,\ldotp\,\ldotp\,\ldotp\,\ldotp\,\ldotp&8,510\,1594\\
\cos\zeta^*\ldotp\,\ldotp\,\ldotp&9,848\,1718_n\\\cline{2-2}
&3,716\,9794_n\\
(1)\ldotp\,\ldotp\,\ldotp\,\ldotp\,\ldotp\,\ldotp\,\ldotp\,\ldotp\,\ldotp\,\ldotp&737
\end{array}\)
\par\vspace{3.7ex}
\(\begin{aligned}
\Delta\varphi&=5212'',584\\
\Delta\zeta&=6279,380\\
\Delta\lambda&=8170,868
\end{aligned}\)
\end{minipage}\hfill
\begin{minipage}[t]{.317\linewidth}\vspace{0pt}\centering
\(\begin{array}{l@{\;}r}
s\ldotp\,\ldotp\,\ldotp\,\ldotp\,\ldotp\,\ldotp&5,358\,6482\\
B\ldotp\,\ldotp\,\ldotp\,\ldotp\,\ldotp\,\ldotp&8,508\,9799\\
\sin\zeta^*\ldotp\,\ldotp\,\ldotp\,\ldotp&9,850\,7902_n\\\cline{2-2}
&3,718\,4183_n\\
\operatorname{tang}\varphi^*\ldotp\,\ldotp&0,079\,4351\\
\cos\varphi^*\ldotp\,\ldotp\,\ldotp&9,806\,1553\\\cline{2-2}
&3,797\,8534_n\\
(2)\ldotp\,\ldotp\,\ldotp\,\ldotp\,\ldotp\,\ldotp\,\ldotp\,\ldotp\,\ldotp\,\ldotp\,\ldotp\,\ldotp&633\\
&3,912\,2630_n\\
(3)\ldotp\,\ldotp\,\ldotp\,\ldotp\,\ldotp\,\ldotp\,\ldotp\,\ldotp\,\ldotp\,\ldotp\,\ldotp\,\ldotp\,\ldotp&52
\end{array}\)
\end{minipage}\hfill
\begin{minipage}[t]{.387\linewidth}\vspace{0pt}\centering
\(\begin{array}{lr@{\quad}lr}
\alpha\ldotp\,\ldotp\,\ldotp\,\ldotp&4,62872&\beta\ldotp\,\ldotp\,\ldotp\,\ldotp&4,62990\\
\delta\zeta^2\ldotp\,\ldotp\,\ldotp&7,59571&\delta\lambda^2\ldotp\,\ldotp\,\ldotp&7,82453\\
\mathrm{I}=&167.66&\mathrm{II}=&284.73\\
2\mathrm{II}=&569.46&&\\
2\mathrm{III}=&465.46&&\\
\gamma\ldotp\,\ldotp\,\ldotp\,\ldotp&1,64955&\delta\ldotp\,\ldotp\,\ldotp\,\ldotp&4,62924\\
ss\ldotp\,\ldotp\,\ldotp&0,71730&\delta\varphi^2\ldotp\,\ldotp&7,43396\\
\mathrm{III}=&232.73&\mathrm{IV}=&115.66
\end{array}\)
\end{minipage}
\par\medskip
[2. Anwendung der Formeln aus Art.~6.]\quad Direkte Methode.
\par\medskip
\noindent\begin{minipage}[t]{.34\linewidth}\vspace{0pt}\centering
\(\begin{array}{lr}
\operatorname{tang}\varphi\ldotp\,\ldotp\,\ldotp\,\ldotp\,\ldotp\,\ldotp&0,068\,3004\\
\sqrt{1-ee}\ldotp\,\ldotp\,\ldotp\,\ldotp\,\ldotp\,\ldotp&-1\,4372\\\cline{2-2}
\operatorname{tang}\psi\ldotp\,\ldotp\,\ldotp\,\ldotp\,\ldotp\,\ldotp&0,066\,8632
\end{array}\)
\par\vspace{7ex}
\(\varphi-\psi=5'37'',206\)
\end{minipage}\hfill
\begin{minipage}[t]{.64\linewidth}\vspace{0pt}\centering
oder genauer: \(\displaystyle\frac{1-\sqrt{1-ee}}{1+\sqrt{1-ee}}\sin(\varphi+\psi)=\sin(\varphi-\psi)\)
\par\medskip
\(\begin{array}{lr}
\displaystyle\frac1\rho.\frac{1-\sqrt{1-ee}}{1+\sqrt{1-ee}}\ldotp\,\ldotp\,\ldotp\,\ldotp\,\ldotp\,\ldotp&2,533\,1319\\[1.3ex]
\sin(\varphi+\psi)\ldotp\,\ldotp\,\ldotp\,\ldotp\,\ldotp\,\ldotp&9,994\,7628\\
&+2\\\cline{2-2}
\varphi-\psi\ldotp\,\ldotp\,\ldotp\,\ldotp\,\ldotp\,\ldotp&2,527\,8949
\end{array}\)
\end{minipage}
\par\endgroup

% Printed88; principalPDF95.
\clearpage
\noindent\makebox[12mm][l]{88}\makebox[\dimexpr\linewidth-24mm\relax][c]{\textsc{Nachlass.}}\makebox[12mm][r]{}\par\medskip
\begingroup\small
\noindent\begin{minipage}[t]{.52\linewidth}\vspace{0pt}\centering
\(\begin{array}{rcl}
\psi&=&49^\circ23'35'',724\\
90^\circ-\psi&=&40\quad36\quad24,276\\
45^\circ-\tfrac12\psi&=&20\quad18\quad12,138\\
\tfrac12S&=&1\quad1\quad38,582\\
45^\circ-\tfrac12\psi-\tfrac12S=\mathrm{I}&=&19^\circ16'33'',556\\
45^\circ-\tfrac12\psi+\tfrac12S=\mathrm{II}&=&21\quad19\quad50,720\\
\tfrac12\zeta&=&112\quad9\quad1,050
\end{array}\)
\end{minipage}\hfill
\begin{minipage}[t]{.46\linewidth}\vspace{0pt}\centering
\(\begin{array}{lr}
1:\sqrt{1-ee}\ldotp\,\ldotp\,\ldotp\,\ldotp\,\ldotp\,\ldotp\,\ldotp\,\ldotp&0,001\,4372.07\\
1:\sqrt{1-ee\sin\varphi^{*2}}\ldotp\,\ldotp\,\ldotp\,\ldotp\,\ldotp\,\ldotp\,\ldotp\,\ldotp&8474.43\\\cline{2-2}
&0,000\,5897.64\\
a\rho\ldotp\,\ldotp\,\ldotp\,\ldotp&1,490\,1727.02\\
&1,489\,5829\\
s\ldotp\,\ldotp\,\ldotp\,\ldotp&5,358\,6482\\\cline{2-2}
S\ldotp\,\ldotp\,\ldotp\,\ldotp&3,869\,0653\\
\multicolumn2c{S=7397'',164}
\end{array}\)
\end{minipage}
\par\medskip
\noindent\begin{minipage}[t]{.65\linewidth}\vspace{0pt}\centering
\(\begin{array}{lr@{\quad}lr}
\sin\mathrm{I}\ldotp\,\ldotp&9,518\,6703&\sin\mathrm{II}\ldotp\,\ldotp&9,560\,8046\\
\sin\tfrac12\zeta\ldotp\,\ldotp&9,966\,7039&\cos\tfrac12\zeta\ldotp\,\ldotp&9,576\,3844_n\\
\cos\mathrm{I}\ldotp\,\ldotp&9,974\,9441&\cos\mathrm{II}\ldotp\,\ldotp&9,969\,1810\\\hline
\sin\tfrac12\zeta\sin\mathrm{I}\ldotp\,\ldotp&9,485\,3742&\sin\tfrac12\zeta\cos\mathrm{I}\ldotp\,\ldotp&9,941\,6480\\
\cos\tfrac12\zeta\sin\mathrm{II}\ldotp\,\ldotp&9,137\,1890_n&\cos\tfrac12\zeta\cos\mathrm{II}\ldotp\,\ldotp&9,545\,5654_n\\
\sin\tfrac12(\zeta'-L)\ldotp\,\ldotp&9,960\,1922&\sin\tfrac12(\zeta'+L)\ldotp\,\ldotp&9,967\,5138
\end{array}\)
\par\vspace{2.7ex}
\(\begin{array}{lr@{\quad}lr}
\sin(45^\circ-\tfrac12\psi')\ldotp\,\ldotp&9,525\,1820&\cos(45^\circ-\tfrac12\psi')\ldotp\,\ldotp&9,974\,1342
\end{array}\)
\par\smallskip
\(\psi'=50^\circ50'31'',424\)
\par\smallskip
\noindent\begin{minipage}[t]{.45\linewidth}\vspace{0pt}\centering
\(\begin{array}{lr}
\operatorname{tang}\psi'\ldotp\,\ldotp\,\ldotp&0,089\,1844\\
\sqrt{1-ee}\ldotp\,\ldotp\,\ldotp\,\ldotp\,\ldotp&+1\,4372\\\cline{2-2}
&0,090\,6216
\end{array}\)
\par\vspace{3ex}
\(\underline{\underline{\varphi'=50^\circ56'5'',536}}\)
\end{minipage}\hfill
\begin{minipage}[t]{.53\linewidth}\vspace{0pt}\centering
\(\begin{array}{lr}
\displaystyle\frac1\rho.\frac{1-\sqrt{1-ee}}{1+\sqrt{1-ee}}\ldotp\,\ldotp&2,533\,1319\\[1.2ex]
\sin(\varphi'+\psi')\ldotp\,\ldotp\,\ldotp&9,990\,7603\\
&+2\\\cline{2-2}
&2,523\,8924
\end{array}\)
\par\smallskip
\(\varphi'-\psi'=334'',112\)
\end{minipage}
\end{minipage}\hfill
\begin{minipage}[t]{.33\linewidth}\vspace{0pt}\centering
\(\begin{array}{rcl}
\tfrac12(\zeta'-L)&=&114^\circ9'31'',723\\
\tfrac12(\zeta'+L)&=&111\quad53\quad9,737\\\cline{3-3}
\mathord{\raisebox{-.2ex}{.}}\,\zeta'&=&226\quad2\quad41,460\\\cline{1-3}\noalign{\vskip .45pt}\cline{1-3}
-L&=&2\quad16\quad21,986\\
&=&8181'',986
\end{array}\)
\par\smallskip
\(\begin{array}{lr}
L\ldotp\,\ldotp\,\ldotp\,\ldotp\,\ldotp&3,912\,8588_n\\[1.3ex]
\displaystyle\sqrt{\frac{1-ee}{1-ee\sin\varphi^{*2}}}\ldotp\,\ldotp\,\ldotp&-5898\\\cline{2-2}
\lambda\ldotp\,\ldotp\,\ldotp\,\ldotp\,\ldotp&3,912\,2690_n
\end{array}\)
\par\smallskip
\(\underline{\underline{-\lambda=8170'',883}}\)
\end{minipage}
\par\medskip
Daraus ergeben sich die Abweichungen von den Ergebnissen der vorherigen Methode
\[
\begin{array}{lr}
\text{beim Azimut}&-0'',020\\
\text{bei der Länge}&+0,015\\
\text{bei der Polhöhe}&+0,022.
\end{array}
\]
\endgroup
\begin{center}[8.]\\[.5ex][Geodätische Übertragung auf der Kugel.]\end{center}
[Für die Kugel gilt näherungsweise, wenn \(R\) den Bogen des Großkreises
geteilt durch den Radius bezeichnet]

% Printed89; principalPDF96.
\clearpage
\noindent\makebox[12mm][l]{}\makebox[\dimexpr\linewidth-24mm\relax][c]{\textsc{Geodätische Linie.}}\makebox[12mm][r]{89}\par\medskip
\[
\begin{aligned}
\Delta\varphi&=R\cos\zeta\frac{\sin\tfrac12R}{\tfrac12R}.
\frac{\tfrac12\Delta\varphi}{\sin\tfrac12\Delta\varphi}\sec\tfrac12\Delta\lambda\\[1ex]
\Delta\zeta&=R\sin\zeta\operatorname{tang}\varphi\frac{\tfrac12\Delta\zeta}{\sin\tfrac12\Delta\zeta}.
\frac{\operatorname{tang}\tfrac12R}{\tfrac12R}\\[1ex]
\Delta\lambda&=R\frac{\sin\zeta}{\cos\varphi}.\frac{\sin\tfrac12R}{\tfrac12R}.
\frac{\tfrac12\Delta\lambda}{\sin\tfrac12\Delta\lambda}
\end{aligned}
\]
\noindent[wobei \(\zeta\) und \(\varphi\) wieder die Mittelwerte der Azimute und Breiten
des Großkreises an den Endpunkten bezeichnen].
\begin{center}[9.]\\[1ex]
[Berechnung der linearen Länge der geodätischen Linie und ihrer Azimute\\
aus den geografischen Koordinaten.]
\end{center}
[Es bezeichne
\begin{center}
\begin{tabular}{r@{\quad}p{.87\linewidth}}
\(r\)&die lineare Länge der geodätischen Linie auf einer Rotationsfläche\\
\(\varphi_1\)&die geografische Breite des Anfangspunkts\\
\(\varphi_1+t\)&die geografische Breite des Endpunkts\\
\(N\)&das südwestliche Azimut des Endpunkts am Anfangspunkt\\
\(M\)&das nordöstliche Azimut des Anfangspunkts am Endpunkt\\
\(\lambda\)&den Längenunterschied.
\end{tabular}
\end{center}
\noindent\(\lambda\) gilt als positiv, wenn der Endpunkt westlich des Anfangspunkts liegt.]
\[
\begin{aligned}
dr&=\surd(ppdt^2+\pi\pi d\lambda^2)\\
p&=a+bt+\tfrac12ctt+\tfrac16dt^3+\ldots\\
\pi&=\alpha+\beta t+\tfrac12\gamma tt+\tfrac16\delta t^3+\ldots.
\end{aligned}
\]
\noindent[Nach den Disqu. gen. c. superf. curv., Art.~22, Band IV, S.~249, gilt:]
\[
\begin{aligned}
\frac{\partial r}{\partial t}&=-p\cos M&
m\frac{\partial N}{\partial\lambda}&=\pi\cos M\\
\frac{\partial r}{\partial\lambda}&=\pi\sin M=\alpha\sin N&
m\frac{\partial N}{\partial t}&=p\sin M.
\end{aligned}
\]
[Aus der Gleichung
\[
\left(\frac1p\frac{\partial rr}{\partial t}\right)^2+
\left(\frac1\pi\frac{\partial rr}{\partial\lambda}\right)^2=4rr
\]
\noindent folgt:]
\par\vfill\noindent IX.\hfill12

% Printed90; principalPDF97.
\clearpage
\noindent\makebox[12mm][l]{90}\makebox[\dimexpr\linewidth-24mm\relax][c]{\textsc{Nachlass.}}\makebox[12mm][r]{}\par\medskip
\begingroup\small
\[
\begin{aligned}
rr={}&aatt+\alpha\alpha\lambda\lambda+abt^3+\alpha\beta t\lambda\lambda\\
&+(\tfrac13ac+\tfrac14bb)t^4+
\left(\tfrac16\frac{b\alpha\beta}{a}+\tfrac13\alpha\gamma\right)tt\lambda\lambda
-\tfrac1{12}\frac{\alpha\alpha\beta\beta}{aa}\lambda^4\ldots
\end{aligned}
\]
\noindent[und hieraus folgt, da
\[
-r\cos M=\frac1{2p}\frac{\partial rr}{\partial t}\quad\text{und}\quad
r\sin M=\frac1{2\pi}\frac{\partial rr}{\partial\lambda}
\]
\noindent gilt:]
\[
\begin{aligned}
-r\cos M&=at+\tfrac12btt+\tfrac12\frac{\alpha\beta}{a}\lambda\lambda+\tfrac16ct^3
+\left(\tfrac13\frac{\alpha\gamma}{a}-\tfrac13\frac{b\alpha\beta}{aa}\right)t\lambda\lambda\ldots\\[1ex]
r\sin M&=\alpha\lambda\qquad\ast\qquad
+\left(\tfrac16\frac{b\beta}{a}-\tfrac16\gamma\right)tt\lambda
-\tfrac16\frac{\alpha\beta\beta}{aa}\lambda^3\ldots.
\end{aligned}
\]
\noindent[Aus den Gleichungen
\[
\begin{aligned}
r\cos N&=-\frac1p r\cos M\frac{\partial(r\cos N)}{\partial t}
+\frac1\pi r\sin M\frac{\partial(r\cos N)}{\partial\lambda}\\[1ex]
r\sin N&=-\frac1p r\cos M\frac{\partial(r\sin N)}{\partial t}
+\frac1\pi r\sin M\frac{\partial(r\sin N)}{\partial\lambda}\\
&=\frac{r\pi}{\alpha}\sin M
\end{aligned}
\]
\noindent ergibt sich weiter:]
\[
\begin{aligned}
-r\cos N&=at+\tfrac12btt-\tfrac12\frac{\alpha\beta}{a}\lambda\lambda+\tfrac16ct^3
+\left(\tfrac16\frac{b\alpha\beta}{aa}-\tfrac16\frac{\alpha\gamma}{a}-\tfrac12\frac{\beta\beta}{a}\right)t\lambda\lambda\ldots\\[1ex]
r\sin N&=\alpha\lambda\qquad+\beta t\lambda\qquad
+\left(\tfrac16\frac{b\beta}{a}+\tfrac13\gamma\right)tt\lambda
-\tfrac16\frac{\alpha\beta\beta}{aa}\lambda^3\ldots
\end{aligned}
\]
[Ferner gilt:]
\[
\begin{aligned}
-\cos M\frac{\partial p}{\partial\lambda}+p\sin M\frac{\partial M}{\partial\lambda}
&=\sin M\frac{\partial\pi}{\partial t}+\pi\cos M\frac{\partial M}{\partial t}\\[1ex]
p\sin M\frac{\partial N}{\partial\lambda}&=\pi\cos M\frac{\partial N}{\partial t},
\end{aligned}
\]
\noindent[Da hier \(\frac{\partial p}{\partial\lambda}=0\) gilt, folgt für
\[N-M=\Delta\]
\noindent die Gleichung:]
\[
\pi r\cos M\frac{\partial\Delta}{\partial t}
-pr\sin M\frac{\partial\Delta}{\partial\lambda}=r\sin M\frac{\partial\pi}{\partial t}
\]
\noindent[oder]
\begingroup\footnotesize
\[
\begin{aligned}
&(\alpha+\beta t+\tfrac12\gamma tt\ldots)
\left(at+\tfrac12btt+\tfrac12\frac{\alpha\beta}{a}\lambda\lambda+\tfrac16ct^3+
\left(\tfrac13\frac{\alpha\gamma}{a}-\tfrac13\frac{b\alpha\beta}{aa}\right)t\lambda\lambda\ldots\right)\frac{\partial\Delta}{\partial t}\\[1ex]
&+(a+bt+\tfrac12ctt\ldots)
\left(\alpha\lambda\qquad\ast\qquad+
\left(\tfrac16\frac{b\beta}{a}-\tfrac16\gamma\right)tt\lambda-\tfrac16\frac{\alpha\beta\beta}{aa}\lambda^3\ldots\right)\frac{\partial\Delta}{\partial\lambda}\\[1ex]
&=-(\beta+\gamma t+\tfrac12\delta tt\ldots)
\left(\alpha\lambda\qquad\ast\qquad+
\left(\tfrac16\frac{b\beta}{a}-\tfrac16\gamma\right)tt\lambda-\tfrac16\frac{\alpha\beta\beta}{aa}\lambda^3\ldots\right)
\end{aligned}
\]
\endgroup\endgroup

% Printed91; principalPDF98.
\clearpage
\noindent\makebox[12mm][l]{}\makebox[\dimexpr\linewidth-24mm\relax][c]{\textsc{Geodätische Linie.}}\makebox[12mm][r]{91}\par\medskip
\begingroup\small
\noindent[oder]
\begingroup\footnotesize
\[
\begin{aligned}
&\left\{1+\frac\beta\alpha t+\left(\tfrac23\frac\gamma\alpha-\tfrac16\frac{b\beta}{a\alpha}\right)tt
+\tfrac16\frac{\beta\beta}{aa}\lambda\lambda\ldots\right\}\\[1ex]
&\quad\left\{at+\tfrac12btt+\tfrac12\frac{\alpha\beta}{a}\lambda\lambda+\tfrac16ct^3+
\left(\tfrac13\frac{\alpha\gamma}{a}-\tfrac13\frac{b\alpha\beta}{aa}\right)t\lambda\lambda\ldots\right\}\frac{\partial\Delta}{\partial t}
+(a+bt+\tfrac12ctt\ldots)\lambda\frac{\partial\Delta}{\partial\lambda}\\[1ex]
&\hfill=-(\beta+\gamma t+\tfrac12\delta tt\ldots)\lambda,
\end{aligned}
\]
\endgroup
\noindent[woraus folgt:]
\[
\begin{aligned}
-\Delta={}&\frac\beta a\lambda+\frac{a\gamma-b\beta}{2aa}t\lambda+
\left(\tfrac16\frac\delta a-\tfrac14\frac{b\gamma}{aa}-\tfrac16\frac{c\beta}{aa}
-\tfrac16\frac{\beta\gamma}{a\alpha}+\tfrac14\frac{bb\beta}{a^3}+\tfrac16\frac{b\beta\beta}{aa\alpha}\right)tt\lambda\\[1ex]
&\hfill-\frac{\alpha\beta(a\gamma-b\beta)}{12a^4}\lambda^3\ldots.
\end{aligned}
\]
\noindent[Setzt man
\[\frac{M+N}{2}=A,\]
\noindent also \(M=A-\tfrac12\Delta\) und \(N=A+\tfrac12\Delta\), so erhält man zunächst:]
\[
\begin{aligned}
r\cos A\sin\tfrac12\Delta&=\tfrac12\beta t\lambda+\tfrac14\gamma tt\lambda\ldots\\[1ex]
-r\sin A\sin\tfrac12\Delta&=\tfrac12\frac{\alpha\beta}{a}\lambda\lambda+
\left(\tfrac14\frac{\alpha\gamma}{a}+\tfrac14\frac{\beta\beta}{a}-\tfrac14\frac{b\alpha\beta}{aa}\right)t\lambda\lambda\ldotp\ldotp:\\[1ex]
-r\cos A\cos\tfrac12\Delta&=at+\tfrac12btt+\tfrac16ct^3+
\left(\tfrac1{12}\frac{\alpha\gamma}{a}-\tfrac14\frac{\beta\beta}{a}-\tfrac1{12}\frac{b\alpha\beta}{aa}\right)t\lambda\lambda\ldots\\[1ex]
r\sin A\cos\tfrac12\Delta&=\alpha\lambda+\tfrac12\beta t\lambda+
\left(\tfrac16\frac{b\beta}{a}+\tfrac1{12}\gamma\right)tt\lambda-\tfrac16\frac{\alpha\beta\beta}{aa}\lambda^3\ldots,
\end{aligned}
\]
\noindent[somit]
\[
\begin{aligned}
-r\cos A&=at+\tfrac12btt+\tfrac16ct^3+
\left(\tfrac1{12}\frac{\alpha\gamma}{a}-\tfrac18\frac{\beta\beta}{a}-\tfrac1{12}\frac{b\alpha\beta}{aa}\right)t\lambda\lambda\ldots\\[1ex]
r\sin A&=\alpha\lambda+\tfrac12\beta t\lambda+
\left(\tfrac16\frac{b\beta}{a}+\tfrac1{12}\gamma\right)tt\lambda-\tfrac1{24}\frac{\alpha\beta\beta}{aa}\lambda^3\ldots.
\end{aligned}
\]
[Drückt man die Koeffizienten \(a,b,c\ldots,\ \alpha,\beta,\gamma\ldots,\) zum Argument
\(\varphi_1\) durch die Werte \(a',b',c'\ldots,\ \alpha',\beta',\gamma'\ldots,\) zum Argument
\(\varphi_1+\tfrac12t\) aus, setzt man also
\(a=a'-\tfrac12b't+\tfrac18c'tt-\tfrac1{48}d't^3\ldots\)
usw., so ergibt sich:]
\[
\begin{aligned}
-r\cos A&=a't+\tfrac1{24}c't^3+
\left(\tfrac1{12}\frac{\alpha'\gamma'}{a'}-\tfrac18\frac{\beta'\beta'}{a'}-\tfrac1{12}\frac{b'\alpha'\beta'}{a'a'}\right)t\lambda\lambda\ldots\\[1ex]
r\sin A&=\alpha'\lambda+
\left(\tfrac16\frac{b'\beta'}{a'}-\tfrac1{24}\gamma'\right)tt\lambda
-\tfrac1{24}\frac{\alpha'\beta'\beta'}{a'a'}\lambda^3\ldots\\[1ex]
-\Delta&=\frac{\beta'}{a'}\lambda+
\left(\tfrac1{24}\frac{\delta'}{a'}-\tfrac1{24}\frac{c'\beta'}{a'a'}-\tfrac16\frac{\beta'\gamma'}{a'\alpha'}
+\tfrac16\frac{b'\beta'\beta'}{a'a'\alpha'}\right)tt\lambda
-\frac{\alpha'\beta'(a'\gamma'-b'\beta')}{12{a'}^4}\lambda^3\ldots.
\end{aligned}
\]
[Für das abgeplattete Rotationsellipsoid gilt, wenn zur Abkürzung
\par\vfill\hfill12*\par
\endgroup
\endgroup
